% chapters_parte2/07_cloud_networking.tex

\chapter{\eh{cloud} Cloud Networking}

Las nubes públicas no son simplemente servidores remotos con Internet rápido. Son
redes enteras diseñadas desde cero para escala masiva, multi-tenancy y automatización
via API. Entender su arquitectura diferencia a un ingeniero que \emph{usa} la nube
de uno que \emph{diseña} sobre ella.

% ─────────────────────────────────────────────
\section{\eh{building-construction} VPC --- Virtual Private Cloud}

\begin{concepto}
Una \textbf{VPC (Virtual Private Cloud)} es una red virtual aislada dentro de la
nube pública. Cada cliente tiene su propia VPC con su propio espacio de direcciones IP,
tablas de routing, y políticas de seguridad. El aislamiento entre VPCs lo garantiza
el hypervisor/SDN del proveedor, no hardware dedicado.

Componentes fundamentales (nomenclatura AWS, análogos en GCP/Azure):
\begin{itemize}
    \item \textbf{CIDR block:} rango IP de la VPC (ej. 10.0.0.0/16).
    \item \textbf{Subnets:} subdivisiones de la VPC, asociadas a una zona de disponibilidad.
    \item \textbf{Route Tables:} controlan a dónde va el tráfico de cada subnet.
    \item \textbf{Internet Gateway (IGW):} permite que subnets públicas accedan a Internet.
    \item \textbf{NAT Gateway:} permite que subnets privadas salgan a Internet sin
          ser directamente accesibles.
    \item \textbf{Security Groups:} firewall stateful por instancia (como iptables).
    \item \textbf{NACLs (Network ACLs):} firewall stateless por subnet.
\end{itemize}
\end{concepto}

\subsection{Arquitectura típica multi-tier}

\begin{center}
\begin{tikzpicture}[
    tier/.style={draw=cyan!50!white, fill=darkcardgray, text=textlight,
                 rectangle, rounded corners=4pt, minimum width=7cm,
                 minimum height=0.9cm, font=\small\bfseries, align=center},
    gw/.style={draw=yellow!60!white, fill=darkcardviolet, text=textlight,
               rectangle, rounded corners=4pt, minimum width=2cm,
               minimum height=0.7cm, font=\tiny, align=center},
    arrow/.style={->, thick, draw=gray!60!white},
]
    \node[gw] (igw)  at (0, 6)   {Internet Gateway};
    \node[tier] (pub)  at (0, 4.5) {Subnet Pública: Load Balancer, Bastion};
    \node[gw] (nat)  at (4, 3)   {NAT GW};
    \node[tier] (app)  at (0, 3)   {Subnet Privada: App Servers, APIs};
    \node[tier] (db)   at (0, 1.5) {Subnet Privada: Bases de Datos, Caches};
    \node[tier] (vpc)  at (0, 0)   {VPC: 10.0.0.0/16};

    \draw[arrow] (igw) -- (pub);
    \draw[arrow] (pub) -- (app);
    \draw[arrow] (app) -- (db);
    \draw[arrow] (app) -- (nat);
    \draw[arrow, dashed] (nat) -- node[font=\tiny, text=yellow!70!white, right] {Internet} (igw);
\end{tikzpicture}
\end{center}

\begin{itemize}
    \item Subnet pública: accesible desde Internet (IGW en route table).
    \item Subnet privada de app: no accesible directamente, sale a Internet vía NAT.
    \item Subnet privada de DB: sin acceso a Internet en ninguna dirección.
\end{itemize}

% ─────────────────────────────────────────────
\section{\eh{link} Conectividad entre VPCs y On-Premise}

\subsection{VPC Peering}

\begin{concepto}
Conexión privada entre dos VPCs (mismo o distinto account, misma o distinta región).
El tráfico viaja por la backbone del proveedor, no por Internet. \textbf{No transitivo:}
si A está peered con B, y B con C, A \textbf{no} puede hablar con C a través de B.
\end{concepto}

\subsection{Transit Gateway / Cloud Router}

\begin{concepto}
Solución hub-and-spoke para conectar muchas VPCs y redes on-premise. En lugar de
$N^2$ peerings, cada VPC conecta al Transit Gateway (AWS) o Cloud Router (GCP).
Sí es transitivo. Permite routing más granular y segmentación de tráfico.
\end{concepto}

\subsection{VPN y Direct Connect / Dedicated Interconnect}

\begin{center}
\begin{tabular}{lll}
\toprule
\textbf{Opción} & \textbf{Características} & \textbf{Cuándo usar} \\
\midrule
VPN (IPsec)       & Sobre Internet, encriptado, rápido de desplegar & Dev/test, fallback \\
Direct Connect    & Enlace físico dedicado, bajo jitter, SLA fuerte & Producción crítica \\
PrivateLink       & Acceso a servicios del proveedor sin Internet     & S3, RDS, APIs internas \\
\bottomrule
\end{tabular}
\end{center}

% ─────────────────────────────────────────────
\section{\eh{bar-chart} Load Balancing en la nube}

\begin{concepto}
Los proveedores cloud ofrecen load balancers gestionados que escalan automáticamente.
La elección depende del tipo de tráfico:
\end{concepto}

\begin{center}
\begin{tabular}{llll}
\toprule
\textbf{Tipo} & \textbf{Capa OSI} & \textbf{AWS} & \textbf{GCP} \\
\midrule
L4 TCP/UDP    & Transporte & Network LB (NLB)    & TCP/UDP LB \\
L7 HTTP/HTTPS & Aplicación & Application LB (ALB) & HTTP(S) LB \\
Global anycast & L7 global & CloudFront + ALB     & Global HTTP LB \\
\bottomrule
\end{tabular}
\end{center}

\begin{ejemplo}
\textbf{ALB (Application Load Balancer)} puede enrutar por:
\begin{itemize}
    \item Host header: \texttt{api.ejemplo.com} $\to$ Target Group A;
          \texttt{admin.ejemplo.com} $\to$ Target Group B.
    \item Path: \texttt{/api/v1} $\to$ servicio legacy; \texttt{/api/v2} $\to$ servicio nuevo.
    \item Header HTTP: \texttt{X-Version: beta} $\to$ 10\% canary deployment.
    \item Peso: 90\% tráfico a stable, 10\% a canary (weighted target groups).
\end{itemize}
\end{ejemplo}

% ─────────────────────────────────────────────
\section{\eh{lock} Seguridad de Red en la Nube}

\subsection{Security Groups vs NACLs}

\begin{center}
\begin{tabular}{lll}
\toprule
\textbf{Característica} & \textbf{Security Groups} & \textbf{NACLs} \\
\midrule
Estado         & Stateful (respuesta automática) & Stateless (reglas explícitas ida y vuelta) \\
Nivel          & Instancia (ENI)                 & Subnet \\
Reglas         & Solo allow                      & Allow y deny \\
Evaluación     & Todas las reglas juntas          & Numeradas, primera que aplica \\
Uso típico     & Firewall de host                & Bloqueo de IPs maliciosas \\
\bottomrule
\end{tabular}
\end{center}

\begin{cuidado}
NACLs son stateless: si permites tráfico entrante al puerto 80, necesitas una regla
de salida para los puertos efímeros (1024--65535) para que las respuestas puedan salir.
Con Security Groups no: al ser stateful, la respuesta se permite automáticamente.
\end{cuidado}

\subsection{AWS PrivateLink / GCP Private Service Connect}

\begin{moderno}
\emoji{rocket} Acceso a servicios (S3, DynamoDB, APIs de terceros) sin que el tráfico
salga a Internet. El servicio se expone como un endpoint en tu VPC con IP privada.
El tráfico nunca cruza un IGW. Mitiga riesgo de exfiltración y elimina cargos de
NAT Gateway para tráfico a servicios del proveedor (S3 puede ser millones al mes).
\end{moderno}

% ─────────────────────────────────────────────
\section{\eh{satellite} CDN y Edge Networking}

\begin{concepto}
Una \textbf{CDN (Content Delivery Network)} distribuye contenido estático (imágenes,
JS, CSS, videos) a \textbf{PoPs (Points of Presence)} cercanos a los usuarios finales.
El objetivo: reducir latencia sirviendo desde el PoP más próximo al usuario, no desde
el origen.
\end{concepto}

\begin{center}
\begin{tabular}{lll}
\toprule
\textbf{Provider} & \textbf{CDN} & \textbf{Diferenciador} \\
\midrule
AWS         & CloudFront     & Integración nativa con S3, ALB, Lambda@Edge \\
GCP         & Cloud CDN      & Backbone privado Google, anycast global \\
Cloudflare  & Cloudflare CDN & DDoS protection, Workers edge compute \\
Akamai      & Akamai CDN     & Mayor red de PoPs del mundo \\
\bottomrule
\end{tabular}
\end{center}

\begin{ejemplo}
\textbf{Latencia con y sin CDN para usuario en CDMX accediendo a servidor en Virginia:}
\begin{itemize}
    \item Sin CDN: CDMX $\to$ Internet $\to$ Virginia. RTT $\approx$ 80ms. Carga HTML+JS+CSS
          requiere múltiples round-trips: $>$500ms.
    \item Con CDN (PoP en CDMX): JS/CSS servidos localmente. RTT al PoP $\approx$ 5ms.
          Solo llamadas a API van a Virginia. Carga percibida: $<$100ms.
\end{itemize}
\end{ejemplo}
